\subsection{Stokes Theory}

% **Role:** calculus as local linearization plus global geometric integration.
%
% Generalized Stokes theorem, Green's theorem, classical Stokes theorem, divergence theorem.

Stokes theory relates exterior differentiation to oriented boundary, with Green's theorem, classical Stokes, and divergence as special cases.

\begin{definition} \label{def:affine-chains-and-boundaries} \hlt{Chains and Boundaries}\\
Let $\mathbf{p}_0,\ldots,\mathbf{p}_k\in\R^n$.
\begin{enumerate}[label=(\roman*)]
\setlength{\itemsep}{0pt}
\item The oriented affine \hlt{$k$-simplex} $[\mathbf{p}_0,\ldots,\mathbf{p}_k]$ is the $k$-surface $\sigma:Q^k\to\R^n$, $\sigma(\mathbf{u})\defeq\mathbf{p}_0+\sum\limits_{i=1}^{k}u_i(\mathbf{p}_i-\mathbf{p}_0)$.
\item For a permutation $\pi$, $[\mathbf{p}_{\pi(0)},\ldots,\mathbf{p}_{\pi(k)}]\defeq\operatorname{sgn}(\pi)[\mathbf{p}_0,\ldots,\mathbf{p}_k]$.
\item An \hlt{affine $k$-chain} is $C=\sum_\ell c_\ell\sigma_\ell$; integration is linear: $\int_C\omega\defeq\sum_\ell c_\ell\int_{\sigma_\ell}\omega$.
\item For $k\geq1$, $\partial_{\mathrm{ch}}[\mathbf{p}_0,\ldots,\mathbf{p}_k]\defeq\sum\limits_{j=0}^{k}(-1)^j[\mathbf{p}_0,\ldots,\mathbf{p}_{j-1},\mathbf{p}_{j+1},\ldots,\mathbf{p}_k]$, extended linearly.\\
If $\Phi=T\circ\sigma$ is a $C^2$ oriented simplex, set $\partial_{\mathrm{ch}}\Phi\defeq T(\partial_{\mathrm{ch}}\sigma)$, extended linearly to chains.
\end{enumerate}
\end{definition}

\begin{theorem} \label{thm:stokes-theorem-differential-forms}\label{thm:forms-manifolds-stokes-main-source-todo} \hlt{Stokes' Theorem}\\
Let $V\subseteq\R^m$ be open, let $\Psi$ be a $C^2$ $k$-chain in $V$, and let $\omega$ be a $C^1$ $(k-1)$-form on $V$. Then $\int_{\Psi}d\omega=\int_{\partial_{\mathrm{ch}}\Psi}\omega$.
\begin{proof}
By linearity, take $\Psi=\Phi=T\circ\sigma$, where $\sigma$ is the standard oriented $k$-simplex.\\
By Theorem \ref{thm:pullback-basic-laws}, $\int_{\Phi}d\omega=\int_{\sigma}d(\omega_T)$ and $\int_{\partial_{\mathrm{ch}}\Phi}\omega=\int_{\partial_{\mathrm{ch}}\sigma}\omega_T$; it remains to prove the formula on $\sigma$.\\
By linearity, reduce to $\lambda=f\,dx_1\wedge\cdots\wedge\widehat{dx_r}\wedge\cdots\wedge dx_k$.\\
For $\partial_{\mathrm{ch}}\sigma=\sum\limits_{j=0}^{k}(-1)^j\sigma_j$, all faces vanish except the two with varying $x_r$ endpoints.\\
The remaining equality is the one-variable FTC in the $x_r$ integral.\\
\hlt{Structural Commentary}: Pullback reduces Stokes to the standard simplex.\\
\hlt{Key Insight}: Boundary signs match the missing wedge factor.
\end{proof}
\end{theorem}

\begin{definition} \label{def:closed-and-exact-forms} \hlt{Closed and Exact Forms}\\
Let $E\subseteq\R^n$ be open and let $\omega$ be a $k$-form on $E$.
\begin{enumerate}[label=(\roman*)]
\setlength{\itemsep}{0pt}
\item $\omega$ is \hlt{exact} if $\omega=d\lambda$ for some $(k-1)$-form $\lambda$ on $E$.
\item $\omega$ is \hlt{closed} if $\omega\in C^1$ and $d\omega=0$.
\end{enumerate}
\end{definition}

\begin{perspective} \label{per:closed-exact-local-global} \hlt{Closed Is Local, Exact Is Global}\\
Exact $\Rightarrow$ closed is just $d^2=0$; the converse depends on the domain.\\
On $\R^2\setminus\{\ze\}$ the standard warning example is $\omega=\frac{-y}{x^2+y^2}\,dx+\frac{x}{x^2+y^2}\,dy$, which is closed but not exact.\\
Holes can obstruct global antiderivatives even when local derivative tests vanish.
\end{perspective}

\begin{theorem} \label{thm:poincare-lemma-convex-rudin} \hlt{Poincare Lemma on Convex Sets}\\
Let $k\geq1$.
\begin{enumerate}[label=(\roman*)]
\setlength{\itemsep}{0pt}
\item If $E\subseteq\R^n$ is open and convex, then every closed $C^1$ $k$-form on $E$ is exact.
\item The same conclusion holds on any open set $C^2$-diffeomorphic to a convex open set.
\end{enumerate}
\begin{proof}
\begin{enumerate}[label=(\roman*)]
\setlength{\itemsep}{0pt}
\item Fix $\mathbf{a}\in E$ and set $H_t(\x)\defeq\mathbf{a}+t(\x-\mathbf{a})$, $0\leq t\leq1$.\\
For a $k$-form $\omega$, define $(K\omega)_{\x}(\mathbf{v}_1,\ldots,\mathbf{v}_{k-1})\defeq\int_0^1t^{k-1}\omega_{H_t(\x)}(\x-\mathbf{a},\mathbf{v}_1,\ldots,\mathbf{v}_{k-1})\,dt$.\\
Differentiating under the integral gives $\omega-H_0^\ast\omega=dK\omega+Kd\omega$.\\
Since $H_0^\ast\omega=0$ for $k\geq1$, $d\omega=0$ implies $\omega=d(K\omega)$.
\item Let $\Phi:C\to E$ be a $C^2$ diffeomorphism from convex open $C$. If $d\omega=0$, then $d(\Phi^\ast\omega)=\Phi^\ast(d\omega)=0$.\\
By (i), $\Phi^\ast\omega=d\eta$ on $C$, hence $\omega=(\Phi^{-1})^\ast d\eta=d((\Phi^{-1})^\ast\eta)$ on $E$.
\end{enumerate}
\hlt{Structural Commentary}: Convex contraction gives a homotopy operator.\\
\hlt{Key Insight}: Closed forms become derivatives along the contraction.
\end{proof}
\end{theorem}

\begin{theorem} \label{thm:vector-analysis-forms-dictionary} \hlt{Vector Analysis Dictionary in $\R^3$}\\
Let $E\subseteq\R^3$ be open and let $\mathbf{F}=F_1\mathbf{e}_1+F_2\mathbf{e}_2+F_3\mathbf{e}_3\in C^1(E,\R^3)$.\\
Set $\lambda_{\mathbf{F}}\defeq F_1\,dx+F_2\,dy+F_3\,dz$ and $\omega_{\mathbf{F}}\defeq F_1\,dy\wedge dz+F_2\,dz\wedge dx+F_3\,dx\wedge dy$.
\begin{enumerate}[label=(\roman*)]
\setlength{\itemsep}{0pt}
\item For $u\in C^1(E)$, $\mathbf{F}=\nabla u$ iff $\lambda_{\mathbf{F}}=du$.
\item $d\lambda_{\mathbf{F}}=\omega_{\nabla\times\mathbf{F}}$; hence $\nabla\times\mathbf{F}=0$ iff $d\lambda_{\mathbf{F}}=0$.
\item For $\mathbf{G}\in C^1(E,\R^3)$, $\mathbf{F}=\nabla\times\mathbf{G}$ iff $\omega_{\mathbf{F}}=d\lambda_{\mathbf{G}}$.
\item $d\omega_{\mathbf{F}}=(\nabla\cdot\mathbf{F})\,dx\wedge dy\wedge dz$; hence $\nabla\cdot\mathbf{F}=0$ iff $d\omega_{\mathbf{F}}=0$.
\item If $E$ is $C^2$-diffeomorphic to a convex open set, then $\nabla\times\mathbf{F}=0\Rightarrow\exists u\in C^1(E):\mathbf{F}=\nabla u$;\\
$\nabla\cdot\mathbf{F}=0\Rightarrow\exists \mathbf{G}\in C^1(E,\R^3):\mathbf{F}=\nabla\times\mathbf{G}$.
\end{enumerate}
\begin{proof}
\begin{enumerate}[label=(\roman*)]
\setlength{\itemsep}{0pt}
\item This is $du=(D_1u)\,dx+(D_2u)\,dy+(D_3u)\,dz$.
\item Direct expansion gives $d\lambda_{\mathbf{F}}=(D_2F_3-D_3F_2)\,dy\wedge dz+(D_3F_1-D_1F_3)\,dz\wedge dx+(D_1F_2-D_2F_1)\,dx\wedge dy$.
\item By (ii), $d\lambda_{\mathbf{G}}=\omega_{\nabla\times\mathbf{G}}$; since $\mathbf{H}\mapsto\omega_{\mathbf{H}}$ is injective, $\omega_{\mathbf{F}}=d\lambda_{\mathbf{G}}$ iff $\mathbf{F}=\nabla\times\mathbf{G}$.
\item Direct expansion gives $d\omega_{\mathbf{F}}=(D_1F_1+D_2F_2+D_3F_3)\,dx\wedge dy\wedge dz$.
\item If $\nabla\times\mathbf{F}=0$, then $\lambda_{\mathbf{F}}$ is closed by (ii), so Poincare gives $\lambda_{\mathbf{F}}=du$, hence $\mathbf{F}=\nabla u$ by (i). If $\nabla\cdot\mathbf{F}=0$, then $\omega_{\mathbf{F}}$ is closed by (iv), so Poincare gives $\omega_{\mathbf{F}}=d\eta$; writing $\eta=\lambda_{\mathbf{G}}$ gives $\mathbf{F}=\nabla\times\mathbf{G}$ by (iii).
\end{enumerate}
\hlt{Structural Commentary}: Exterior derivative is grad--curl--div.\\
\hlt{Key Insight}: Vector identities are $d^2=0$ in $\R^3$.
\end{proof}
\end{theorem}

\begin{theorem} \label{thm:green-stokes-divergence-special-cases} \hlt{Green, Classical Stokes, and Divergence}\\
With the indicated positive orientations:
\begin{enumerate}[label=(\roman*)]
\setlength{\itemsep}{0pt}
\item For $\Omega\subseteq\R^2$ with positively oriented $\operatorname{bd}\Omega$ and $\alpha,\beta\in C^1$ near $\Omega$, $\int_{\operatorname{bd}\Omega}(\alpha\,dx+\beta\,dy)=\int_{\Omega}(D_1\beta-D_2\alpha)\,dA$.
\item If $\Phi$ is a $C^2$ oriented $2$-surface in $\R^3$ and $\mathbf{F}\in C^1$ near $\Phi$, then $\int_{\partial_{\mathrm{ch}}\Phi}\mathbf{F}\cdot\mathbf{t}\,ds=\int_{\Phi}(\nabla\times\mathbf{F})\cdot\mathbf{n}\,dA$.
\item If $\Omega\subseteq\R^3$ has positively oriented boundary and $\mathbf{F}\in C^1$ near $\Omega$, then $\int_{\operatorname{bd}\Omega}\mathbf{F}\cdot\mathbf{n}\,dA=\int_{\Omega}(\nabla\cdot\mathbf{F})\,dV$.
\end{enumerate}
\begin{proof}
Use the form Stokes identity $\int_{\Psi}d\eta=\int_{\partial_{\mathrm{ch}}\Psi}\eta$.
\begin{enumerate}[label=(\roman*)]
\setlength{\itemsep}{0pt}
\item For $\lambda\defeq\alpha\,dx+\beta\,dy$, $d\lambda=(D_1\beta-D_2\alpha)\,dx\wedge dy$; hence $\int_{\operatorname{bd}\Omega}\lambda=\int_{\Omega}d\lambda$ gives Green's formula.
\item For $\lambda\defeq\lambda_{\mathbf{F}}$, $d\lambda=\omega_{\nabla\times\mathbf{F}}$; hence $\int_{\partial_{\mathrm{ch}}\Phi}\lambda=\int_{\Phi}d\lambda$ gives the classical Stokes formula.
\item For $\omega\defeq\omega_{\mathbf{F}}$, $d\omega=(\nabla\cdot\mathbf{F})\,dx\wedge dy\wedge dz$; hence $\int_{\operatorname{bd}\Omega}\omega=\int_{\Omega}d\omega$ gives the divergence formula.
\end{enumerate}
\hlt{Structural Commentary}: One form identity yields the three vector formulas.\\
\hlt{Key Insight}: Boundary plus $d$ is invariant; notation changes.
\end{proof}
\end{theorem}

\begin{method} \label{meth:forms-manifolds-proof-patterns-source-todo}\label{meth:differential-forms-proof-patterns} \hlt{Differential-Forms Proof Patterns}
\begin{enumerate}[label=(\roman*)]
\setlength{\itemsep}{0pt}
\item \hlt{Signs}: compute wedge signs by determinant swaps and permutation parity.
\item \hlt{Standard presentation}: remove repeated indices and sort the rest.
\item \hlt{Pullback}: substitute $dt_i$ for $dy_i$ and apply the chain rule.
\item \hlt{Stokes}: pull back to the standard simplex, prove the coordinate case, extend linearly.
\item \hlt{Closed versus exact}: use $d^2=0$ for necessity; use convexity or $C^2$-equivalence for primitives.
\item \hlt{Vector analysis}: translate fields to forms, apply Stokes, translate back.
\end{enumerate}
\hlt{Key Insight}: Forms package signs, orientation, and parameter changes.
\end{method}
